\section{Introduction}
\label{sec:introduction}

Submarine cables carry nearly all intercontinental Internet traffic and connect
users to remote services~\cite{telegeography_submarine_map}. A cable or coastal
corridor shared by several routes can expose those routes to the same physical
failure. Evaluating path diversity therefore requires knowing whether routes
that appear distinct at the network layer also occupy distinct physical
directions.

Traceroute records routers and Autonomous System (AS) transitions, so network-
layer measurements count distinct hops or AS paths as distinct alternatives.
It does not identify the submarine infrastructure that could carry those
paths. Two paths that appear independent in traceroute may leave a region
through the same coast and rely on the same small set of physical directions.
Whether network-layer diversity corresponds to physical-corridor diversity is
an empirical question.

Cross-layer systems make this comparison possible. Internet Atlas, InterTubes,
and iGDB connect observed Internet links to long-haul infrastructure;
Nautilus associates IP links with candidate submarine cables; Calypso maps
routes using cable layouts, network relationships, and inland connectivity;
and Xaminer applies cross-layer maps to infrastructure risk analysis
~\cite{durairajan2013atlas,durairajan2015intertubes,anderson2022igdb,
ramanathan2023nautilus,wang2025calypso,ramanathan2024xaminer}. These systems
produce mappings for individual links or paths. We use such mappings to ask a
population-level question: does concentration at the network layer track
concentration across feasible physical corridors?

We present \emph{CLASP}, a Cross-Layer Audit of Service Paths. CLASP normalizes
each traceroute into visible, geolocated hop-pair segments. For every country--
measurement population, ordered endpoint-AS transitions form the network
distribution and feasible pairs of landing regions form the corridor
distribution. Both distributions contain the same candidate-bearing segments.
Each segment contributes one unit of observation mass, allocated among its
candidate corridors by the projection score. We call the Spearman rank
association between network and corridor concentration across a reported group
of units \emph{Layer Agreement}.

We apply CLASP to 490,911 valid IPv4 traceroutes from 18 public RIPE Atlas
measurements. The corpus covers all 13 DNS Roots, Wikipedia, Reddit, Netflix
Assets, and two dynamic multi-target measurements used as topology references.
The aligned one-hour snapshot yields 3,291,063 atomic segments, of which
2,432,559 are mappable and 124,350 support at least one feasible inter-region
corridor. The audit contains 370 country--measurement units: 222 DNS Root, 53
application, and 95 topology-reference units.

The observed Layer Agreement differs between geographic groups. Under
projection-score weighting, it is 0.632 among 67 island and archipelagic
service-facing units and 0.001 among 202 coastal-mainland units. The group
difference remains in the same direction under equal-share and Top-1
allocation. These coefficients describe rank associations within the observed
snapshot; the analysis does not establish that geography causes the
difference.

The service-facing population is also heterogeneous. Of 275 units, 81 (29.5\%)
have a broad network distribution and a concentrated corridor distribution,
while 158 (57.5\%) remain broad at both layers. Other units are concentrated at
both layers or expand after projection. A single contraction rate therefore
does not represent all country--service populations.

This paper makes three contributions. First, CLASP constructs paired network
and corridor distributions from identical traceroute segments. Second, we
measure a geographic difference in the observed association between the two
layers and report it under three existing candidate-allocation rules. Third,
we provide a reproducible audit structure that records measurement IDs,
configuration values, mapping states, and pipeline counts, allowing the
reported distributions and their limitations to be inspected.
